% !TeX root = ../tcc.tex
\chapter*[Referências]{Referências}
\addcontentsline{toc}{chapter}{Referências}

FERNANDES, Jacks Renan Neves; TELES, Ariel Soares; FERNANDES, Thayaná
Ribeiro Silva; LIMA, Lucas Daniel Batista; BALHARA, Surjeet; GUPTA,
Nishu; TEIXEIRA, Silmar. Artificial Intelligence on Diagnostic Aid of
Leprosy: A Systematic Literature Review. \emph{Journal of Clinical
Medicine}, v. 13, n.~1, artigo 180, 2024. DOI:
https://doi.org/10.3390/jcm13010180. \textbf{Nota editorial:} o PDF
informa publicação online em 28/12/2023, mas sua referência e volume são
de 2024; adotou-se o ano bibliográfico de 2024.

JAIKISHORE, Chathura N.; UDUTALAPALLY, Venkanna; DAS, Debanjan. AI
Driven Edge Device for Screening Skin Lesion and Its Severity in
Peripheral Communities. In: \emph{2021 IEEE 18th India Council
International Conference (INDICON)}. IEEE, 2021. p.~1--6. DOI:
https://doi.org/10.1109/INDICON52576.2021.9691666.

YOTSU, Rie R.; DING, Zhengming; HAMM, Jihun; BLANTON, Ronald E. Deep
learning for AI-based diagnosis of skin-related neglected tropical
diseases: A pilot study. \emph{PLOS Neglected Tropical Diseases}, v. 17,
n.~8, e0011230, 2023. DOI: https://doi.org/10.1371/journal.pntd.0011230.

SANCHEZ, Karen; HINOJOSA, Carlos; MIELES, Olinto; ZHAO, Chen; GHANEM,
Bernard; ARGUELLO, Henry. CO2Wounds-V2: Extended Chronic Wounds Dataset
from Leprosy Patients. In: \emph{2024 IEEE International Conference on
Image Processing (ICIP)}. IEEE, 2024. p.~69--75. DOI:
https://doi.org/10.1109/ICIP51287.2024.10647641.

EITA, Mohammad A.; RIZK, Hamada. A trustworthy and explainable deep
learning framework for skin lesion detection in smart dermatology.
\emph{Engineering Applications of Artificial Intelligence}, v. 159,
artigo 111594, 2025. DOI:
https://doi.org/10.1016/j.engappai.2025.111594.

COCO. \emph{Detection Evaluation}. Documentação técnica do benchmark.
Disponível em: https://cocodataset.org/dataset/detection-eval.htm.
Acesso em: 15 out. 2026. Fonte complementar para interpretação de IoU e
AP/mAP em caixas e máscaras; não é estudo incluído na revisão de
hanseníase.

